\begin{proof}[Proof of \cref{obj:lem:block-energy-representer}]
Fix \(\beta\ge1\) and \(p\in(0,1)\).  For every integer \(r\ge1\) write
\[
e_r(p):=\frac{\Delta_r(p)^2}{[p(1-p)]^r},
\qquad
m_r(p):=\frac{|\Delta_r(p)|}{[p(1-p)]^r}(1+p)^r ,
\]
and define the three constants of the statement by
\[
c_1:=\min\bigl\{e_r(p):1\le r\le\beta,\ \Delta_r(p)\ne0\bigr\},
\qquad
c_2:=2^\beta\sum_{r=1}^{\beta}e_r(p),
\qquad
H:=c_1+\frac{2^\beta\sum_{r=1}^{\beta}m_r(p)}{c_1}.
\]
Since \(\Delta_1(p)=(1-p)-(-p)=1\ne0\), the index set defining \(c_1\) contains \(r=1\) and is therefore a nonempty finite set; on it \(\Delta_r(p)^2>0\) and \([p(1-p)]^r>0\), so \(c_1\) is a minimum of finitely many strictly positive numbers and \(c_1>0\).  Each \(m_r(p)\ge0\) for the same reason, so \(H\ge c_1>0\). % lean: blockLowerConst_spec, energyCoeff, rawMassCoeff, blockRawMassConst

\textbf{Step 1 (exposed orders).}
Fix \(d\ge1\) and write \(k_\star:=k_\star(d,\beta,p)\).  Since \(\beta\ge1\) and \(d\ge1\) we have \(\bar\beta_d=\min\{\beta,d\}\ge1\), and \(\Delta_1(p)=1\ne0\) shows that the set \(\{r\in\{1,\ldots,\bar\beta_d\}:\Delta_r(p)\ne0\}\) of \cref{obj:def:exposed-order} is nonempty.  Hence \(k_\star\) is the maximum of that set, so
\[
1\le k_\star\le\bar\beta_d,
\qquad
\Delta_{k_\star}(p)\ne0,
\qquad\text{and}\qquad
r\le k_\star
\ \text{ for every } r\in\{1,\ldots,\bar\beta_d\}\ \text{with}\ \Delta_r(p)\ne0 .
\]
In particular \(k_\star\le\beta\) and \(k_\star\le d\). % lean: kStar_mem_exposedOrder, eligibleOrder_le_kStar

\textbf{Step 2 (binomial domination at fixed order).}
For all integers \(r\le k\le d\) with \(k\le\beta\),
\[
\binom d r\le 2^\beta\binom d k .
\]
Indeed, counting the pairs \(R\subseteq K\subseteq\{1,\ldots,d\}\) with \(|R|=r\) and \(|K|=k\) in two ways gives \(\binom d r\binom{d-r}{k-r}=\binom d k\binom k r\); the left factor \(\binom{d-r}{k-r}\ge1\) and \(\binom k r\le 2^k\le2^\beta\). % lean: choose_le_pow_mul_choose

\textbf{Step 3 (energy scale).}
By \cref{obj:def:block-score-energy},
\[
A_d=\sum_{r=1}^{\bar\beta_d}\binom d r e_r(p),
\]
a sum of nonnegative terms.  Discarding all summands except \(r=k_\star\), which is legitimate by Step 1,
\[
A_d\ \ge\ \binom{d}{k_\star}e_{k_\star}(p)\ \ge\ c_1\binom{d}{k_\star},
\]
the second inequality because \(1\le k_\star\le\beta\) and \(\Delta_{k_\star}(p)\ne0\), so \(e_{k_\star}(p)\) is one of the numbers over which \(c_1\) is the minimum.  For the upper bound fix \(r\in\{1,\ldots,\bar\beta_d\}\).  If \(\Delta_r(p)=0\) then \(e_r(p)=0\) and both sides of the next display vanish; if \(\Delta_r(p)\ne0\) then \(r\le k_\star\le\min\{\beta,d\}\) by Step 1 and Step 2 applies.  Either way
\[
\binom d r e_r(p)\le 2^\beta\binom{d}{k_\star}e_r(p).
\]
Summing over \(r\in\{1,\ldots,\bar\beta_d\}\) and then enlarging the index range from \(\{1,\ldots,\bar\beta_d\}\) to \(\{1,\ldots,\beta\}\), which only adds nonnegative terms since \(\bar\beta_d\le\beta\),
\[
A_d\le 2^\beta\binom{d}{k_\star}\sum_{r=1}^{\beta}e_r(p)=c_2\binom{d}{k_\star}.
\]
This is the energy-scale claim. % lean: blockEnergy_uniform_compare

\textbf{Step 4 (\(c_1\le c_2\)).}
Apply Step 3 at \(d=1\) and set \(K:=\binom{1}{k_\star(1,\beta,p)}\).  Step 1 gives \(k_\star(1,\beta,p)\le\bar\beta_1\le1\), so \(K\ge1>0\), and
\[
c_1K\le A_1\le c_2K .
\]
Dividing by \(K>0\) gives \(c_1\le c_2\). % lean: blockEnergy_uniform_compare, kStar_mem_exposedOrder

\textbf{Step 5 (design moments).}
Fix \(d\ge1\) and let \(D\) be a finite design on \(\{0,1\}^d\) satisfying the common-\(p\) product Bernoulli condition of \cref{obj:ass:bernoulli-design}.  Then \(D\) assigns to each \(z\in\{0,1\}^d\) the mass \(\prod_{j=1}^d p^{z_j}(1-p)^{1-z_j}>0\), so for every \(f:\{0,1\}^d\to\mathbb R\),
\[
\mathbb E_D[f(Z)^2]=0
\qquad\text{if and only if}\qquad
f(z)=0\ \text{ for every } z\in\{0,1\}^d .
\]
% lean: blockDesign_sq_eq_zero_iff
Moreover, factorizing the expectation over the independent coordinates, for every nonempty \(S\subseteq\{1,\ldots,d\}\) and every \(T\subseteq\{1,\ldots,d\}\),
\[
\mathbb E_D\!\left[\prod_{j\in S}(Z_j-p)\prod_{j\in T}Z_j\right]
=
\begin{cases}
[p(1-p)]^{|S|}\,p^{|T|-|S|}, & S\subseteq T,\\
0, & \text{otherwise},
\end{cases}
\]
\[
\mathbb E_D\!\left[\prod_{j\in S}(Z_j-p)\prod_{j\in T}(Z_j-p)\right]
=
\begin{cases}
[p(1-p)]^{|S|}, & S=T,\\
0, & \text{otherwise}.
\end{cases}
\]
In the first identity a coordinate in \(S\cap T\) contributes \(\mathbb E[(Z_j-p)Z_j]=p(1-p)\), a coordinate of \(T\) outside \(S\) contributes \(\mathbb E[Z_j]=p\), a coordinate of \(S\) outside \(T\) contributes \(\mathbb E[Z_j-p]=0\), and all remaining coordinates contribute \(1\); in the second identity a coordinate in \(S\cap T\) contributes \(\mathbb E[(Z_j-p)^2]=p(1-p)\) and a coordinate of the symmetric difference contributes \(0\). % lean: E_centeredMonomial_mul_raw, E_centeredMonomial_mul

\textbf{Step 6 (the score represents \(L_d\)).}
For every integer \(t\ge1\),
\[
\sum_{r=1}^{t}\binom t r\Delta_r(p)\,p^{t-r}=1,
\]
because \(\sum_{r=0}^{t}\binom t r(1-p)^rp^{t-r}=\bigl[(1-p)+p\bigr]^t=1\) and \(\sum_{r=0}^{t}\binom t r(-p)^rp^{t-r}=(-p+p)^t=0\), while the two \(r=0\) terms cancel since \(\Delta_0(p)=0\). % lean: blockContrast_binomial
Let \(T\subseteq\{1,\ldots,d\}\) be nonempty with \(|T|\le\bar\beta_d\).  Inserting the definition of \(g_d\) from \cref{obj:def:block-score-energy} into the first moment identity of Step 5, and using that exactly \(\binom{|T|}{r}\) of the sets \(S\) with \(|S|=r\) satisfy \(S\subseteq T\),
\[
\mathbb E_D\!\left[g_d(Z)\prod_{j\in T}Z_j\right]
=\sum_{r=1}^{\bar\beta_d}\frac{\Delta_r(p)}{[p(1-p)]^r}\binom{|T|}{r}[p(1-p)]^{r}p^{|T|-r}
=\sum_{r=1}^{\bar\beta_d}\binom{|T|}{r}\Delta_r(p)\,p^{|T|-r}
=1,
\]
where the last equality is the previous display with \(t=|T|\ge1\): the summands with \(|T|<r\le\bar\beta_d\) vanish because \(\binom{|T|}{r}=0\), so the sum runs effectively over \(1\le r\le|T|\). % lean: blockScore_raw_moment
Taking \(T=\varnothing\) in the same moment identity, where no nonempty \(S\) satisfies \(S\subseteq T\),
\[
\mathbb E_D[g_d(Z)]=0 .
\]
% lean: blockScore_mean_zero
Now \(L_d\bigl(z\mapsto\prod_{j\in T}z_j\bigr)=1\) for nonempty \(T\) and \(L_d(1)=0\), so the two displays state precisely that \(\mathbb E_D[g_d(Z)f(Z)]=L_d(f)\) whenever \(f\) is a raw monomial of degree at most \(\bar\beta_d\).  Both sides are linear in \(f\), and \cref{obj:def:perturbation-program} defines \(\mathcal P_d\) as the span of these raw monomials; hence for every \(f\in\mathcal P_d\),
\[
\mathbb E_D[g_d(Z)f(Z)]=L_d(f).
\]
% lean: blockScore_represents

\textbf{Step 7 (energy identities and positivity).}
Evaluating \(g_d\) at the all-treated and at the all-control assignment, each of the \(\binom d r\) sets \(S\) with \(|S|=r\) contributes \((1-p)^r\) and \((-p)^r\) respectively, so
\[
L_d(g_d)=\sum_{r=1}^{\bar\beta_d}\frac{\Delta_r(p)}{[p(1-p)]^r}\binom d r\bigl[(1-p)^r-(-p)^r\bigr]
=\sum_{r=1}^{\bar\beta_d}\binom d r e_r(p)=A_d .
\]
% lean: blockScore_contrast
By the second moment identity of Step 5 the centered monomials appearing in \(g_d\) are pairwise orthogonal and the one indexed by \(S\) has second moment \([p(1-p)]^{|S|}\), so
\[
\mathbb E_D[g_d(Z)^2]
=\sum_{r=1}^{\bar\beta_d}\binom d r\frac{\Delta_r(p)^2}{[p(1-p)]^{2r}}[p(1-p)]^{r}
=A_d .
\]
% lean: blockScore_sq_expectation
Keeping only the \(r=1\) summand of \(A_d\) and using \(\Delta_1(p)=1\),
\[
A_d\ \ge\ \binom d 1\frac{\Delta_1(p)^2}{p(1-p)}=\frac{d}{p(1-p)}>0 .
\]
% lean: blockEnergy_pos
Since \(h_d=g_d/A_d\) with \(A_d>0\), dividing the two preceding identities by \(A_d\) and by \(A_d^2\) gives
\[
L_d(h_d)=1,
\qquad
\mathbb E_D[h_d(Z)^2]=A_d^{-1}.
\]
% lean: blockRepresenter_contrast_energy

\textbf{Step 8 (the score lies in \(\mathcal P_d\)).}
Expanding the product coordinate by coordinate, for every \(S\subseteq\{1,\ldots,d\}\) and every \(z\),
\[
\prod_{j\in S}(z_j-p)=\sum_{T\subseteq S}(-p)^{|S|-|T|}\prod_{j\in T}z_j .
\]
% lean: centeredMonomial_raw_expansion
Every raw monomial occurring on the right has degree \(|T|\le|S|\le\bar\beta_d\) when \(|S|\le\bar\beta_d\), hence lies in \(\mathcal P_d\).  As \(g_d\) is a finite linear combination of centered monomials with \(|S|\le\bar\beta_d\), it follows that \(g_d\in\mathcal P_d\) and therefore \(h_d=g_d/A_d\in\mathcal P_d\).  Together with \(L_d(h_d)=1\) from Step 7, \(h_d\) is feasible for the program of \cref{obj:def:perturbation-program}. % lean: blockScore_mem_polySpace, blockRepresenter_perturbFeasible

\textbf{Step 9 (perturbation optimality).}
Let \(h\in\mathcal P_d\) satisfy \(L_d(h)=1\).  Step 6 applied to \(f=h\) gives \(\mathbb E_D[g_d(Z)h(Z)]=1\), and Step 7 gives \(\mathbb E_D[g_d(Z)^2]=A_d\).  Expanding the square,
\[
\mathbb E_D\!\left[(A_dh(Z)-g_d(Z))^2\right]
=A_d^2\,\mathbb E_D[h(Z)^2]-2A_d\,\mathbb E_D[g_d(Z)h(Z)]+\mathbb E_D[g_d(Z)^2]
=A_d^2\,\mathbb E_D[h(Z)^2]-A_d .
\]
The left-hand side is nonnegative and \(A_d>0\), so \(A_d^2\,\mathbb E_D[h(Z)^2]\ge A_d\), that is,
\[
A_d^{-1}\le\mathbb E_D[h(Z)^2].
\]
Equality holds if and only if the displayed square has expectation zero, which by the faithfulness statement of Step 5 happens if and only if \(A_dh(z)=g_d(z)\) for every \(z\), that is, if and only if \(h=h_d\).  Conversely \(h=h_d\) attains the value \(A_d^{-1}\) by Step 7. % lean: perturbFeasible_energy_unique, blockDesign_sq_eq_zero_iff

\textbf{Step 10 (weight optimality).}
Let \(w\) be feasible for \cref{obj:def:weight-program} under \(D\), that is, \(\mathbb E_D[w(Z)]=0\) and \(\mathbb E_D\bigl[w(Z)\prod_{j\in S}Z_j\bigr]=1\) for every nonempty \(S\subseteq\{1,\ldots,d\}\) with \(|S|\le\bar\beta_d\).  Since \(L_d(1)=0\) and \(L_d\bigl(z\mapsto\prod_{j\in S}z_j\bigr)=1\) for nonempty \(S\), these constraints say that \(\mathbb E_D[w(Z)f(Z)]=L_d(f)\) on every raw monomial \(f\) of degree at most \(\bar\beta_d\); both sides are linear in \(f\), and \cref{obj:def:perturbation-program} defines \(\mathcal P_d\) as the span of those raw monomials, so for every \(f\in\mathcal P_d\),
\[
\mathbb E_D[w(Z)f(Z)]=L_d(f).
\]
% lean: weightFeasibleAt_represents
Taking \(f=g_d\), which lies in \(\mathcal P_d\) by Step 8, and using \(L_d(g_d)=A_d\) and \(\mathbb E_D[g_d(Z)^2]=A_d\) from Step 7,
\[
\mathbb E_D\!\left[(w(Z)-g_d(Z))^2\right]
=\mathbb E_D[w(Z)^2]-2\,\mathbb E_D[w(Z)g_d(Z)]+\mathbb E_D[g_d(Z)^2]
=\mathbb E_D[w(Z)^2]-A_d .
\]
Nonnegativity of the left-hand side gives
\[
A_d\le\mathbb E_D[w(Z)^2],
\]
and equality forces the displayed square to have expectation zero, hence \(w=g_d\) by the faithfulness statement of Step 5.  Conversely, the two identities of Step 6 are exactly the feasibility constraints for \(w=g_d\), and \(\mathbb E_D[g_d(Z)^2]=A_d\), so \(g_d\) is feasible and attains the value \(A_d\). % lean: blockScore_mem_polySpace, blockScore_weightFeasibleAt, weightFeasibleAt_energy_unique, blockDesign_sq_eq_zero_iff

\textbf{Step 11 (program values).}
Both programs are infima of sets of expectations of squares, hence of sets of nonnegative reals, and both sets are nonempty: \(h_d\) is feasible for the perturbation program with value \(A_d^{-1}\) by Steps 7 and 8, and \(g_d\) is feasible for the weight program with value \(A_d\) by Step 10.  Therefore
\[
\mathsf{PerturbProg}_{d,\beta,p}\le A_d^{-1},
\qquad
\mathsf{WeightProg}_{d,\beta,p}\le A_d ,
\]
while Steps 9 and 10 bound every value in the two feasible sets below by \(A_d^{-1}\) and by \(A_d\) respectively, giving the reverse inequalities.  Hence
\[
\mathsf{PerturbProg}_{d,\beta,p}=A_d^{-1},
\qquad
\mathsf{WeightProg}_{d,\beta,p}=A_d .
\]
% lean: blockPrograms_exact

\textbf{Step 12 (coefficient envelope).}
Expanding each centered monomial of \(g_d\) by the identity of Step 8, collecting the coefficient of \(\prod_{j\in T}z_j\), and dividing by \(A_d\), the raw coefficients of \(h_d\) in \cref{obj:def:block-score-energy} are
\[
h_{d,T}
=A_d^{-1}\sum_{r=1}^{\bar\beta_d}\ \
\sum_{\substack{S\subseteq\{1,\ldots,d\},\ |S|=r\\ T\subseteq S}}
\frac{\Delta_r(p)}{[p(1-p)]^r}(-p)^{\,r-|T|}.
\]
% lean: blockRepresenter_raw_expansion
Applying the triangle inequality inside each \(h_{d,T}\) and then exchanging the order of summation, so that each \(S\) with \(|S|=r\) is paired with the subsets \(T\subseteq S\),
\[
\sum_{T\subseteq\{1,\ldots,d\}}|h_{d,T}|
\ \le\
A_d^{-1}\sum_{r=1}^{\bar\beta_d}\ \sum_{\substack{S\subseteq\{1,\ldots,d\}\\ |S|=r}}\ \sum_{T\subseteq S}\frac{|\Delta_r(p)|}{[p(1-p)]^r}\,p^{\,r-|T|}
\ =\
A_d^{-1}\sum_{r=1}^{\bar\beta_d}\binom d r m_r(p),
\]
using \(\sum_{T\subseteq S}p^{|S|-|T|}=(1+p)^{|S|}=(1+p)^r\) and the fact that there are \(\binom d r\) sets \(S\) with \(|S|=r\). % lean: blockRawCoef_mass_estimate
Exactly as in Step 3 — orders with \(\Delta_r(p)=0\) contribute zero, orders with \(\Delta_r(p)\ne0\) satisfy \(r\le k_\star\) and hence \(\binom d r\le2^\beta\binom{d}{k_\star}\) by Step 2, and the index range may then be enlarged to \(\{1,\ldots,\beta\}\) —
\[
\sum_{r=1}^{\bar\beta_d}\binom d r m_r(p)\le 2^\beta\binom{d}{k_\star}\sum_{r=1}^{\beta}m_r(p).
\]
Combining this with the lower energy bound \(c_1\binom{d}{k_\star}\le A_d\) of Step 3, with \(A_d>0\) and \(\binom{d}{k_\star}>0\),
\[
\sum_{T\subseteq\{1,\ldots,d\}}|h_{d,T}|
\ \le\
\frac{2^\beta\binom{d}{k_\star}\sum_{r=1}^{\beta}m_r(p)}{A_d}
\ \le\
\frac{2^\beta\sum_{r=1}^{\beta}m_r(p)}{c_1}
\ \le\ H,
\]
the last step because \(H\) exceeds that quotient by \(c_1>0\). % lean: blockRawCoef_mass_uniform
\end{proof}
